% ===== PRÉAMBULE COMMUN — À COPIER TEL QUEL EN TÊTE DE CHAQUE FICHIER .tex =====
% Ne modifier QUE :
%   - les deux lignes \fancyhead[L] et \fancyhead[R]
%   - le bloc titre \begin{center}...\end{center} juste après \begin{document}
\documentclass[11pt,a4paper]{article}
\usepackage[utf8]{inputenc}
\usepackage[T1]{fontenc}
\usepackage{lmodern}
\usepackage[french]{babel}
\usepackage{amsmath,amssymb,mathtools}
\usepackage{icomma}
\usepackage{xcolor}
\usepackage{tikz}
\usepackage{pgfplots}
\usepackage[most]{tcolorbox}
\usepackage{enumitem}
\usepackage{array,booktabs,colortbl}
\usepackage[a4paper,margin=2.2cm,headheight=26pt,headsep=10pt]{geometry}
\usepackage{fancyhdr}
\usepackage[hidelinks]{hyperref}
\usetikzlibrary{arrows.meta,calc,angles,quotes,positioning,decorations.markings,intersections,backgrounds}
\pgfplotsset{compat=1.18}

\definecolor{primary}{RGB}{226,74,88}
\definecolor{primarydark}{RGB}{150,28,42}
\definecolor{methcol}{RGB}{40,90,150}
\definecolor{excol}{RGB}{0,135,90}
\definecolor{defcol}{RGB}{0,114,178}
\definecolor{attncol}{RGB}{200,40,55}
\definecolor{thmcol}{RGB}{0,140,120}
\definecolor{courbe}{RGB}{32,110,200}
\definecolor{courbeder}{RGB}{210,70,40}
\definecolor{tang}{RGB}{0,150,90}
\definecolor{grille}{RGB}{200,205,215}
\definecolor{plus}{RGB}{200,60,60}
\definecolor{moins}{RGB}{30,90,200}

\tcbset{boxbase/.style={enhanced,breakable,boxrule=0.9pt,arc=2.5pt,
  left=3mm,right=3mm,top=2.3mm,bottom=2.3mm,fonttitle=\bfseries,coltitle=white}}
\newtcolorbox{cle}[1][]{boxbase,colback=defcol!6,colframe=defcol,
  title={\textsc{À retenir}\ifx\relax#1\relax\relax\else\textmd{ : \itshape #1}\fi}}
\newtcolorbox{methode}[1][]{boxbase,colback=methcol!7,colframe=methcol,sharp corners,
  title={\textsc{Méthode}\ifx\relax#1\relax\relax\else\textmd{ --- \itshape #1}\fi}}
\newtcolorbox{exemplebox}[1][]{boxbase,colback=excol!5,colframe=excol,
  title={\textsc{Exemple}\ifx\relax#1\relax\relax\else\textmd{ : \itshape #1}\fi}}
\newtcolorbox{piege}[1][]{boxbase,colback=attncol!6,colframe=attncol,
  title={\textsc{Attention}\ifx\relax#1\relax\relax\else\textmd{ : \itshape #1}\fi}}
\newtcolorbox{exo}[1][]{boxbase,colback=primary!4,colframe=primary,
  title={\textsc{Exercice}\ifx\relax#1\relax\relax\else\textmd{~#1}\fi}}
\newtcolorbox{corr}[1][]{boxbase,colback=thmcol!5,colframe=thmcol,
  title={\textsc{Corrigé}\ifx\relax#1\relax\relax\else\textmd{~#1}\fi}}

\newcommand{\R}{\mathbb{R}}
\newcommand{\N}{\mathbb{N}}
\newcommand{\ds}{\displaystyle}
\newcommand{\tg}{\operatorname{tg}}
\newcommand{\cotg}{\operatorname{cotg}}
\newcommand{\arctg}{\operatorname{arctg}}
\newcommand{\arccotg}{\operatorname{arccotg}}
\newcommand{\limh}{\lim_{h\to 0}}

\pagestyle{fancy}\fancyhf{}
\fancyhead[L]{\small\textcolor{primarydark}{\textbf{Thème 3} — Fonctions composées}}
\fancyhead[R]{\small\textcolor{primarydark}{\textsc{Corrigé — Moyen}}}
\fancyfoot[C]{\small\thepage}
\renewcommand{\headrulewidth}{0.8pt}

\begin{document}

\begin{center}
{\LARGE\bfseries\color{primarydark} Arcs, exponentielles, produits — corrigé}\\[2pt]
{\color{primary}\rule{0.5\textwidth}{1.2pt}}\\[4pt]
{\small\itshape Niveau moyen}
\end{center}
\vspace{2mm}

\begin{corr}[1 --- $y=\arctg(3x)$ puis $y=\arcsin(x^2)$]
$\bullet$ $U=3x$, $U'=3$, formule $(\arctg U)'=\dfrac{U'}{1+U^2}$ :
\[
y'=\frac{3}{1+(3x)^2}=\frac{3}{1+9x^2}.
\]
$\bullet$ $U=x^2$, $U'=2x$, formule $(\arcsin U)'=\dfrac{U'}{\sqrt{1-U^2}}$ :
\[
y'=\frac{2x}{\sqrt{1-(x^2)^2}}=\frac{2x}{\sqrt{1-x^4}}.
\]
\end{corr}

\begin{corr}[2 --- $y=\log_{10}(x)$ puis $y=2^{x}$]
$\bullet$ Formule $(\log_a U)'=\dfrac{U'}{U\ln a}$ avec $U=x$, $U'=1$, $a=10$ :
\[
y'=\frac{1}{x\ln 10}.
\]
$\bullet$ Formule $(a^U)'=U'\,a^U\ln a$ avec $U=x$, $U'=1$, $a=2$ :
\[
y'=2^{x}\ln 2.
\]
\end{corr}

\begin{corr}[3 --- $y=\tg(2x)$ puis $y=\cotg(x)$]
$\bullet$ $U=2x$, $U'=2$, formule $(\tg U)'=\dfrac{U'}{\cos^2 U}$ :
\[
y'=\frac{2}{\cos^2(2x)}.
\]
$\bullet$ $U=x$, $U'=1$, formule $(\cotg U)'=-\dfrac{U'}{\sin^2 U}$ :
\[
y'=-\frac{1}{\sin^2 x}.
\]
\end{corr}

\begin{corr}[4 --- $y=x^2 e^{x}$]
Règle du produit $(uv)'=u'v+uv'$ avec $u=x^2$ ($u'=2x$) et $v=e^x$ ($v'=e^x$) :
\[
y'=2x\,e^x+x^2 e^x = x(2+x)e^x.
\]
\end{corr}

\begin{corr}[5 --- $y=e^{x}\sin(x)$]
Règle du produit avec $u=e^x$ ($u'=e^x$) et $v=\sin x$ ($v'=\cos x$) :
\[
y'=e^x\sin x + e^x\cos x = e^x(\sin x+\cos x).
\]
\end{corr}

\end{document}
